\chapter{Einleitung}
\label{chap:einleitung}
Mit dem Gedankenexperiment einer Katze, die zugleich lebendig und tot ist, veranschaulichte Erwin Schrödinger 1935~\cite{schroedinger1935} eine der verblüffendsten Konsequenzen der Quantenmechanik. Phänomene der Quantenphysik wie die Superposition, bei der sich ein Teilchen gleichzeitig in mehreren Zuständen befinden kann, oder die Verschränkung, durch die Teilchen über beliebige Distanzen hinweg miteinander korreliert bleiben, widersprechen der Alltagserfahrung grundlegend. Genau von diesen Effekten erhofft man sich im Bereich des Computings einen fundamentalen Fortschritt~\cite{Taherdoost2025}. Das entsprechende Feld nennt sich Quantencomputing und nutzt die Superposition, um Systeme zu bauen, deren kleinste Recheneinheiten nicht nur die klassischen Zustände 0 und 1 annehmen können, sondern beide zugleich.

Die Forschung an Quantencomputern hat in den vergangenen Jahren erheblich an Aufschwung gewonnen. Sowohl große Technologiekonzerne wie IBM, Google und Microsoft als auch spezialisierte Unternehmen wie IonQ oder D-Wave und zahlreiche Universitäten investieren beträchtliche Mittel in die Entwicklung leistungsfähiger Quantensysteme und Quantensoftware~\cite{gibney2019quantum,bova2021commercial}. Auch staatliche Stellen stufen die Entwicklung von Quantencomputern als strategisch relevant ein. Der Bund stellte seit 2017 rund eine Milliarde Euro allein für das Quantencomputing bereit. Über sämtliche Quantentechnologien hinweg summieren sich die öffentlichen Zuwendungen von Bund, Ländern und EU auf über 1,9~Milliarden Euro~\cite{bundestag2025quantentechnologien}. Dies zeigt, dass es sich um mehr als ein theoretisches Konzept handelt. Vielversprechende Anwendungsfelder liegen unter anderem in der Simulation von Molekülen für die Arzneimittelentwicklung~\cite{santagati2024drug}, in Optimierungsproblemen aus Finanzwesen und Logistik, in der Kryptographie~\cite{bova2021commercial} sowie im maschinellen Lernen~\cite{biamonte2017quantum}. Dabei handelt es sich um Problemklassen, bei denen klassische Rechner an praktische Grenzen stoßen.

Zum Programmieren dieser Quantencomputer werden größtenteils Python-basierte Frameworks genutzt~\cite{fingerhuth2018opensource}, mit denen sich sogenannte Schaltkreise und Gatter definieren lassen, die die Quantenbits manipulieren. In der täglichen Arbeit am Fraunhofer-Institut für Produktionstechnik und Automatisierung wird nicht nur eines dieser Frameworks eingesetzt, da sie jeweils unterschiedliche Vor- und Nachteile aufweisen. Diese Arbeit beschränkt sich auf die Frameworks \textit{PennyLane} und \textit{Qiskit}, die zu den populärsten zählen~\cite{wang_advantages_2025} und auch bei Fraunhofer häufig verwendet werden. Beide nutzen zur Ausführung unterschiedliche Backends und verwenden verschiedene Notationen. Daraus ergibt sich das Problem, dass sich dasselbe Programm nicht auf beliebigen Backends ausführen lässt, sondern zunächst auf die jeweilige Notation umgeschrieben werden muss.

Am Fraunhofer IPA wurde hierfür ein Executor entwickelt, eine Software, die mehrere Frameworks bündelt und Schaltkreise an das jeweils gewünschte Backend weiterleitet. In der bestehenden Version werden Gatter und Qubits über einen Wrapper definiert, der \textit{Qiskit}-Objekte in die Objekte der einzelnen Frameworks übersetzt. Dadurch ist der Executor eng an \textit{Qiskit} gebunden, was der eigentlichen Aufgabe des Executors, nämlich der Framework-Unabhängigkeit, entgegensteht.

Ziel dieser Arbeit ist es, diese Bindung durch eine framework-unabhängige Abstraktionsschicht innerhalb des Executors aufzulösen. Über diese lassen sich Schaltkreise und Gatter einheitlich definieren, die anschließend automatisch an das gewünschte Backend weitergeleitet werden. Den Nutzerinnen und Nutzern bleibt dadurch das wiederholte Neudefinieren beim Wechsel des Frameworks erspart, und sie können innerhalb desselben Programms verschiedene Frameworks verwenden. Über diese Vereinheitlichung hinaus entstehen weitere Vorteile. Da ein Schaltkreis nur einmal beschrieben und unverändert auf unterschiedlichen Backends ausgeführt wird, beruhen Backend-Vergleiche auf einer kontrollierten Problemformulierung. Zudem lassen sich die einzelnen Frameworks unabhängig voneinander installieren, und Änderungen an deren Schnittstellen bleiben auf eine zentrale Schicht begrenzt, sodass ein einmal geschriebenes Programm über Versionswechsel hinweg lauffähig bleibt.

Das Vorgehen und die Struktur der weiteren Kapitel orientieren sich an drei Forschungsfragen. Zunächst wird untersucht, wie sich eine solche Abstraktionsschicht entwerfen und implementieren lässt (FF1). Darauf aufbauend wird bewertet, welche Vorteile die Abstraktionsschicht gegenüber der direkten Nutzung der einzelnen Frameworks bietet (FF2). Abschließend wird analysiert, welcher Zeit-Overhead durch den Executor im Vergleich zur direkten Nutzung von \textit{PennyLane} und \textit{Qiskit} entsteht (FF3).